\documentclass{ximera}
\title{The Gradient Vector CW}

\newcommand{\pskip}{\vskip 0.1 in}

\begin{document}
\begin{abstract}
The gradient vector and its properties
\end{abstract}
\maketitle

\section{Working in Two Dimensions}

\begin{question} \label{Q3f3r}
The illustration below shows the level curve (blue) of a temperature function through the point $P$. 

Suppose at $P$ the temperature increases at the rate of $0.2^\circ$/meter in the direction of the vector ${\bf v}$.

Use common sense, not a formula, to approximate the direction in which the temperature increases at the greatest rate from $P$. Use sliders $\theta_1$ and $s$ in Lines 1 and 2 to approximate this greatest rate of change. Include units.

\begin{onlineOnly}
    \begin{center}
\desmos{rasqej1ixl}{900}{600}
\end{center}
\end{onlineOnly}

Desmos activity available at

\href{https://www.desmos.com/calculator/rasqej1ixl}{163: Gradient 1}
\end{question}



\begin{question} \label{Qpdf3r40}

At point $P$ below, the temperature in increases at the respective rates of $0.2^\circ$C/m and $0.4^\circ$C/m in the direction of the vectors ${\bf v}$ and ${\bf w}$. Use common sense and the sliders in Lines 1-5 to approximate the direction from $P$ in which the temperature increases at the greatest rate. Then approximate this rate of change.

\begin{onlineOnly}
    \begin{center}
\desmos{4zibriggsa}{900}{600}
\end{center}
\end{onlineOnly}

Desmos activity available at

\href{https://www.desmos.com/calculator/4zibriggsa}{163: Gradient 2}


\end{question}




\begin{question}  \label{Qd4r4356}
The figure below shows two level curves of the differentiable function $z=f(x,y)$ and the gradient of that function at point $C$. Sketch the gradient at point $D$. Explain your reasoning.

\pdfOnly{
Access Geogebra interactives through the online version of this text at
 
\href{https://www.geogebra.org/classic/zau4grcm}.
}
 
\begin{onlineOnly}
    \begin{center}
\geogebra{zau4grcm}{900}{600}
\end{center}
\end{onlineOnly}

Geogebra activity available at

\href{https://www.geogebra.org/classic/zau4grcm}{163: Gradient Exercise}
\end{question}


\begin{question} \label{Q34g4vbty}
The function 
\[
 \sigma = f(x,y)
\]
gives the charge density (in $\mu C/m^2$, ie. micro-Coulombs/squar meter) at the point with $(x,y)$ (measured in meters) on the surface of a thin conducting plate.

Two level curves of the function $f$ are shown below.

\pskip

\pdfOnly{
Access Geogebra interactives through the online version of this text at
 
\href{https://www.geogebra.org/classic/xnqmhyvd}.
}
 
\begin{onlineOnly}
    \begin{center}
\geogebra{xnqmhyvd}{900}{600}
\end{center}
\end{onlineOnly}

Geogebra activity available at

\href{https://www.geogebra.org/classic/xnqmhyvd}{163: Gradient Exercise 2}


Suppose point $C$ has coordinates $(2,5)$, that
\[
    f(2,5) = 4,
\]
and that
\[
    \Big|  \langle \frac{\partial \sigma}{\partial x}, \frac{\partial \sigma}{\partial y} \rangle \Big|_{(x,y) =(2,5)} = 30
\]

\pskip

The gradient vector at $C$ is shown in the figure above. Consecutive marks along the line $CD$ are at intervals of length $1$cm. Answer the following questions \emph{without} imposing a coordinate system. Work with the vectors themselves, not with their components.

(a) What are the units of the above gradient? Explain its meaning.

(b) Approximate the charge density at $D$.

(c) Approximate the charge density at $F$.



\end{question}




\begin{question}  \label{Qweeregtt5}
A light at the point $(a,b,c)$ illuminates part of the paraboloid
\[
  k z = x^2 + y^2.
\]
Here $k>0$ is a constant.

\pdfOnly{
Access Geogebra interactives through the online version of this text at
 
\href{https://www.geogebra.org/classic/nfkajvkp}.
}
 
\begin{onlineOnly}
    \begin{center}
\geogebra{nfkajvkp}{900}{600}
\end{center}
\end{onlineOnly}

Geogebra activity available at

\href{https://www.geogebra.org/classic/nfkajvkp}{163: Terminator}



(a) Find an algebraic inequality that gives necessary and sufficient conditions for the point $(x,y,z)$ on the paraboloid to be illuminated.

(b) Parameterize the curve on the surface (the terminator) that forms the boundary separating the illuminated and dark parts of the paraboloid. 

(c) Compute the area of the illuminated part of the paraboloid.
\end{question}


\end{document}